% ===== §5 =====
\section{The Second Impossibility: The Relationality of Value and the Absent Quantifier}
\label{p25:sec:relationality}

\noindent
Suppose the first impossibility set aside, and grant, only to press on, that a function could be written over such values as do admit of terms. The second condition of the anatomy then comes into view: that this function be exogenous and fixed, given before the process of pursuit and unchanged by it, a determinate target rather than a moving one. This section denies that condition. But it denies it in a particular way, and the way matters as much as the denial, because a weak version of the denial changes nothing and the economist has already absorbed it. The weak version says that people's preferences are unstable, that they want one thing today and another tomorrow, and that the objective therefore drifts. To this the economist answers, rightly, that instability is just more structure to model: let the utility function carry a time index, or let today's choice include a term for its effect on tomorrow's tastes, and the drift is absorbed into a larger fixed problem. The paper's denial is not the weak one and does not meet this answer, because it does not locate the trouble in the instability of preferences at all. It locates it deeper, in the ontology of the value itself, and argues that value in intimate relation is \emph{relational}, generated within the relation and existing only there, so that there is no standpoint outside the relation from which its magnitude could be read off and fixed. The objective cannot be exogenous because there is no exogenous position for it to be given from; and this is not a fact about wandering tastes that a time index could catch, but a fact about where relational value is, which no reindexing reaches.

\subsection*{Value as relational rather than intrasubjective}

The utility function is a function of one subject. It encodes \emph{his} preferences, the values as they stand for \emph{him}, and its arguments are states of the world evaluated from his single standpoint; and even the theory of games, which sets several such functions interacting, leaves each the private possession of its owner, given before the interaction and merely played out through it. The whole construction presupposes that value is intrasubjective, lodged inside a subject and carried by him to the situation as a piece of his equipment. It is exactly this presupposition that the relational ontology of the series denies, and the denial is not a mood but a claim with consequences. The value at stake in an intimate relation, the meaning that a private word carries, the weight a kept silence has come to hold, the worth of the bond itself, is not a property either partner keeps inside and brings; it is generated between them, in the relating, and has no existence apart from the relation that bears it \parencite{paperxxiv,paperxiii}. This is not a metaphor. A private language means nothing to either speaker considered alone, because its meaning is the sediment of a shared history and is stored in the relation and not in either head; a rite that is precious to the two is not a datum in either partner's separate ledger, because its preciousness is not a fact about either of them but a fact about what has formed between them. Value here has the standing that trust was shown to have in the paper on the non-binding vow: a property of a relational structure, a feature of the topology of the bond, and not an attribute of any node within it \parencite{paperxiii}.

\begin{casebox}{the worth that is not in either of them}
Two people have, over years, made a certain walk their own: a particular route, taken on particular evenings, that has gathered to itself the whole texture of their life together, so that to walk it is for each of them an act dense with meaning. Ask what the walk is worth, and to whom, and the question comes apart in the hand. It is not worth much to either of them \emph{separately}: taken alone, one of them walking the route is a person on a road, and the meaning that floods the shared walk is simply absent, having no bearer. Nor is its worth the sum of two private worths, for there are no two private worths to sum; the meaning does not exist in either of them to be added. It exists between them, in the relation the walk both expresses and sustains, and it exists nowhere else. There is, then, no ledger in which the walk's value is an entry, because a ledger is kept by a subject and totals what has worth \emph{for him}, and the walk's worth is not of that kind. To insist on entering it one would have to pick a standpoint, his or hers, and record from there a private shadow of a value that is not private, thereby recording something other than the value one meant to capture.
\end{casebox}

\subsection*{The quantifier that quantification requires, and its absence}

From the relationality of value a consequence follows that is fatal to the exogeneity condition, and it is worth drawing out slowly, because it joins this impossibility to the deepest commitment of the series. To \emph{quantify} a value, to assign it the magnitude a term in $J$ demands, is to occupy a definite position: the position of one who surveys the field of value from outside and above, sees the candidates and their worths spread before him, and reads each onto a common scale. Quantification is silently a view from nowhere within the valued field, a God's-eye operation, and it presupposes a standpoint exterior to the values it measures. For value that is intrasubjective this presupposition is met without difficulty, at least in principle, since the subject stands over against the world he evaluates as its external assessor; his preferences are inside him and the world is outside, and he reads the one onto the other. But for value that is \emph{relational}, generated within the relation and existing only there, the presupposition cannot be met at all. Whoever would quantify such value must stand either within the relation or outside it, and neither will serve. Standing within, he is a party to the very generation he would assess, with no exterior purchase, no vantage from which the relational value lies spread before him as an object; he is inside the thing he would survey. Standing outside, he finds nothing to survey, because the value exists only within and is not present to any external position. There is no third place, no exterior and omniscient vantage from which the worth of what two people generate between them could be read onto a scale. And this is precisely the standpoint that quantification requires: the standpoint of a big Other who would hold the true measure of all relational value. The whole of this series, following the mother paper, has argued that this Other does not exist, that there is no metalanguage and no position from which the meanings by which people live could be assessed from without \parencite{wan2024grb,paperxiii}. To assign relational value a determinate magnitude is therefore to presuppose the one standpoint the series has shown to be empty.

\begin{claim}[The vacant quantifier]
Value in intimate relation is \emph{relational}: generated within the relation and existing only there, a property of the relational structure and not of any subject within it. To quantify a value is to assign it a magnitude from a standpoint exterior to it. For relational value no such exterior standpoint exists, since a position within the relation has no external purchase on it and a position outside the relation has nothing before it to assess; the standpoint quantification requires is that of a big Other holding the measure of all relational value, and this position is empty. The objective of an optimization over intimate relation therefore cannot be exogenous, having no external vantage from which to be given, nor fixed, being generated in the same motion that would pursue it. The endogeneity of the objective is not a contingency of unstable preference but the ontological consequence of the relationality of value.
\end{claim}

\subsection*{The defeat of the economist's reply}

This is the point at which the economist's reply is not merely met but defeated, and the defeat should be stated exactly, because the reply is resilient and returns in new dress if left half-answered. The reply was: your objective is misspecified or unstable; specify it better, index it to time, and the trouble dissolves. Against the first impossibility this failed because certain values admit of no term at all. Against the second it fails for a further and independent reason: that even for values which might admit of terms, there is no standpoint from which the terms could be assigned their magnitudes, the values being relational and the assigning position vacant. The reply rests, at bottom, on a picture in which the true worths are already out there, lodged somewhere and waiting to be read by a careful enough observer, so that better specification is a matter of reading them more accurately. That picture is exactly what the relationality of value refuses. Relational value is not lying in wait, outside the relation, to be discovered by a sufficiently patient economist; it is generated inside the relation, in the act, and holds still for no exterior gaze. The objective is not a target the searcher has specified too crudely. It is generated by the searching, revised in the relating, and held fixed by no standpoint that could serve as its source; and an $\argmax$ over an objective that is produced in the same motion that would maximize it, and that no external position holds still, does not name a hard problem awaiting a better economist. It names no problem, for there is nothing of the form the operation requires for the operation to be applied.
